\chapter{Mission Requirements}
\label{ch:mission_requirements}
This chapter defines the mission-level requirements derived from the LINKU technical objectives. These requirements describe the capabilities that the mission must provide to demonstrate tether deployment, observe tethered-system behavior, and support comparison with pre-flight analyses and ground-test results.

\begin{table}[htbp]
\centering
\caption{LINKU mission-level requirements.}
\label{tab:mission_requirements}
\small
\resizebox{\linewidth}{!}{%
\begin{tabular}{p{2.0cm}p{12.8cm}}
\toprule
\textbf{ID} & \textbf{Mission requirement} \\
\midrule
MR-001 &
The LINKU mission shall demonstrate the in-orbit deployment and observation of a two-satellite tethered system composed of SAT-B and SAT-C. \\

MR-002 &
The mission shall deploy a non-conductive tether between SAT-B and SAT-C toward a nominal length of \(100~\mathrm{m}\). \\

MR-003 &
The mission shall generate time-correlated data to characterize SAT-B/SAT-C separation, tether release, tension build-up, possible slack/taut transitions, residual oscillations, and subsequent tethered-system motion. \\

MR-004 &
The mission shall provide data products suitable for comparison with pre-flight tether-dynamics simulations and ground-test results. \\

MR-005 &
The mission shall support the assessment of passive gravity-gradient-influenced behavior after tether deployment. \\
\bottomrule
\end{tabular}%
}
\end{table}